
Code generation with large language models has achieved notable performance on benchmarks such as HumanEval and MBPP. Current systems typically employ execution feedback to iteratively refine generated code until tests pass. Static analysis tools, which provide type safety and compositional guarantees without execution cost, have received less attention in the context of LLM feedback orchestration.

This work addresses the question: does the order and combination of feedback sources affect computational efficiency in LLM code generation? Specifically, we compare cascade routing (static analysis first, then conditional execution) against aggregation (both sources simultaneously).

The research originated from constraints encountered in a prior investigation that attempted comprehensive vulnerability detection with template-based SMT translation. That effort required extensive infrastructure (112K-sample dataset, 42 CWE vulnerability classes, multi-hour training) exceeding proof-of-concept validation scope. The current study adopts a narrower focus: validating feedback routing mechanisms on a smaller, controlled task set.

We test three mechanism hypotheses: (1) static analysis provides substantial early error detection with zero execution cost, (2) sequential single-source feedback presentation reduces iteration count via attention economy, and (3) conditional execution gating maintains token efficiency within acceptable overhead. Hypothesis 2 was not validated due to runtime errors during experiment execution.
